\documentclass[11pt,a4paper]{article}

% Kodowanie i język
\usepackage[T1]{fontenc}
\usepackage[utf8]{inputenc}
\usepackage[polish]{babel}
\usepackage{lmodern}

% Układ strony
\usepackage[a4paper,hmargin=2.2cm,vmargin=2cm]{geometry}
\usepackage{microtype}
\usepackage{booktabs,tabularx,array}
\usepackage{enumitem}
\usepackage{xcolor}
\usepackage{titlesec}
\usepackage{tikz}
\usetikzlibrary{positioning,fit,shapes.multipart,arrows.meta}
\usepackage[hidelinks]{hyperref}

\definecolor{akcent}{HTML}{23476A}

% Akapity i odstępy
\setlength{\parindent}{0pt}
\setlength{\parskip}{3pt}
\linespread{1.05}
\renewcommand{\arraystretch}{1.12}

% Listy
\setlist[itemize]{leftmargin=1.5em,topsep=2pt,itemsep=1pt,parsep=0pt}
\setlist[enumerate]{leftmargin=1.6em,topsep=2pt,itemsep=1pt,parsep=0pt}

% Nagłówki sekcji: kolor akcentu, linia pod spodem
\titleformat{\section}
  {\large\bfseries\color{akcent}}{\thesection.}{0.5em}{}[{\color{akcent!40}\titlerule}]
\titlespacing*{\section}{0pt}{10pt}{4pt}
\titleformat{\subsection}{\normalsize\bfseries\color{akcent}}{\thesubsection.}{0.5em}{}
\titlespacing*{\subsection}{0pt}{6pt}{2pt}

% Znacznik miejsc do uzupełnienia
\newcommand{\douzup}[1]{\textcolor{red!75!black}{\textbf{[DO UZUPEŁNIENIA: #1]}}}

% Style rysunków
\tikzset{
  encja/.style={rectangle split, rectangle split parts=2, draw=akcent,
    rectangle split part fill={akcent,white}, text width=3.3cm, anchor=north,
    font=\scriptsize\ttfamily, align=left, inner sep=2.5pt, rounded corners=1.5pt},
  replika/.style={encja, rectangle split part fill={akcent!55,white}},
  blok/.style={draw=akcent, fill=akcent!8, rounded corners=3pt, align=center,
    font=\small, inner sep=5pt, minimum width=4.4cm},
  rel/.style={draw=akcent!70, thick},
  krt/.style={font=\scriptsize\sffamily, text=akcent},
}
\newcommand{\pk}{\hfill{\sffamily\tiny\color{akcent}PK}}
\newcommand{\fk}{\hfill{\sffamily\tiny\color{gray}FK}}
\newcommand{\pkfk}{\hfill{\sffamily\tiny\color{akcent}PK\color{gray}, FK}}
\newcommand{\naglowek}[1]{\textcolor{white}{\normalfont\bfseries\footnotesize #1}}

\begin{document}

% ============================================
% STRONA TYTUŁOWA
% ============================================
\begin{titlepage}
\centering
\vspace*{1.5cm}
{\large\bfseries POLITECHNIKA WROCŁAWSKA\par}
\vspace{0.3cm}
{\normalsize Kierunek: Informatyczne Systemy Automatyki, studia II stopnia\par}

\vspace{4cm}
{\color{akcent}\rule{0.6\textwidth}{0.8pt}\par}
\vspace{0.6cm}
{\Huge\bfseries\color{akcent} Rozproszone bazy danych\par}
\vspace{0.6cm}
{\LARGE Aplikacja bazodanowa do obsługi sieci sklepów\par}
\vspace{0.4cm}
{\large Sprawozdanie z projektu\par}
\vspace{0.6cm}
{\color{akcent}\rule{0.6\textwidth}{0.8pt}\par}

\vfill
\begin{tabular}{@{}ll@{}}
\textbf{Autorzy:}    & Natalia Polańska (272567) \\
                     & Michał Pawlik (272566) \\[4pt]
\textbf{Prowadzący:} & \douzup{prowadzący, termin zajęć} \\
\end{tabular}

\vspace{1.5cm}
{\small Wrocław, listopad 2026\par}
\end{titlepage}

% ============================================
\section{Cel i zakres zadania}
Celem projektu jest rozproszona baza danych dla sieci dwóch sklepów wraz
z aplikacją dla pracowników. Wrocław korzysta z bazy Oracle, a Warszawa
z SQL Servera. Serwery działają na dwóch komputerach połączonych siecią Tailscale.
Każdy sklep ma w swojej bazie własne stany, ceny, zamówienia i rezerwacje.
Wspólny katalog produktów prowadzi Wrocław, a Warszawa ma jego kopię.

% ============================================
\section{Funkcjonalny opis działania systemu}
Pracownik każdego sklepu może:
\begin{itemize}
  \item sprawdzić ceny i dostępność produktu w obu sklepach,
  \item sprzedać towar z własnego magazynu,
  \item zarezerwować towar dla klienta, także w drugim sklepie, oraz wydać lub anulować rezerwację,
  \item przyjąć towar, zmienić cenę i przekazać towar do drugiego sklepu,
  \item wygenerować raport dzienny sprzedaży całej sieci.
\end{itemize}
Katalog produktów edytuje się tylko we Wrocławiu; Warszawa pobiera jego kopię.
Nie można sprzedać ani zarezerwować więcej sztuk, niż jest dostępnych.
Gdy jeden serwer nie działa, drugi sklep pracuje dalej.

% ============================================
\section{Przyjęte założenia}
\begin{itemize}
  \item \textbf{Klient:} aplikacja webowa w Pythonie (Streamlit), wspólna dla obu
        sklepów. Łączy się z każdą bazą osobno.
  \item \textbf{Bazy:} Oracle AI Database 26ai Free (Wrocław) i SQL Server 2022 Developer
        (Warszawa). Edycja Developer ma wszystkie funkcje edycji Enterprise.
  \item \textbf{Fragmentacja pozioma:} stany, zamówienia i rezerwacje są przechowywane
        w bazie swojego sklepu. Pozycje zamówień leżą razem z zamówieniem
        (fragmentacja pochodna).
  \item \textbf{Podział danych produktu:} dane stałe (nazwa, kategoria) są w katalogu
        wspólnym dla sieci, a dane zmienne (cena, liczba sztuk) w tabeli
        \texttt{stany\_magazynowe} każdego sklepu. Dzięki temu ten sam produkt może mieć
        w każdym sklepie inną cenę. Taki układ (tabele \texttt{kategorie}
        i \texttt{typy\_produktow} zamiast jednej tabeli \texttt{produkty}) przyjęliśmy
        po uwagach prowadzącego do konspektu.
  \item \textbf{Replikacja:} katalog jest kopiowany z Oracle do SQL Servera co 5 minut
        albo na żądanie. Katalog można zmieniać tylko w Oracle, a kopia w SQL Serverze
        jest dla aplikacji tylko do odczytu. Jedno miejsce zmian wyklucza konflikty,
        np. dwa różne produkty o tym samym identyfikatorze.
  \item \textbf{Linked Server:} SQL Server odczytuje dane z Oracle. Służy do replikacji,
        widoków globalnych (dane obu sklepów w jednym widoku) i raportu dziennego.
  \item \textbf{Transakcje:} każda operacja jest lokalną transakcją w jednej bazie.
        Przekazanie towaru zmienia dwie bazy, więc w razie błędu jest cofane
        operacją kompensującą.
\end{itemize}

% ============================================
\section{Opis implementacji baz danych}

\subsection{Struktura systemu}
\begin{center}
\begin{tikzpicture}
  \node[blok, fill=akcent!18] (app) {\textbf{Aplikacja Python (Streamlit)}};
  \node[blok, below left=1.3cm and -0.4cm of app] (ora)
    {\textbf{Oracle -- Wrocław}\\ \footnotesize katalog główny, stany,\\ \footnotesize zamówienia, rezerwacje};
  \node[blok, below right=1.3cm and -0.4cm of app] (mss)
    {\textbf{SQL Server -- Warszawa}\\ \footnotesize kopia katalogu, stany, zamówienia,\\ \footnotesize rezerwacje, raporty};
  \draw[rel, {Latex}-{Latex}] (app.south west) ++(0.6,0) -- (ora.north);
  \draw[rel, {Latex}-{Latex}] (app.south east) ++(-0.6,0) -- (mss.north);
  \draw[rel, dashed, -{Latex}] (mss.west) -- (ora.east)
    node[midway, above, krt] {Linked Server};
  \node[draw=gray!60, dashed, rounded corners=4pt, inner sep=6pt,
        fit=(ora)(mss), label={[font=\scriptsize\sffamily, text=gray]below:sieć Tailscale}] {};
\end{tikzpicture}
\end{center}

\begin{center}
\begin{tikzpicture}
  \node[replika] (ka) at (0,0) {\naglowek{kategorie}
    \nodepart{two} id\_kategorii \pk\\ nazwa};
  \node[replika] (ty) at (4.2,0) {\naglowek{typy\_produktow}
    \nodepart{two} id\_produktu \pk\\ nazwa\\ id\_kategorii \fk\\ opis\\ aktywny};
  \node[encja] (st) at (8.4,0) {\naglowek{stany\_magazynowe}
    \nodepart{two} id\_produktu \pkfk\\ sklep\\ cena\\ ilosc\_dostepna\\ ilosc\_zarezerwowana\\ data\_modyfikacji};
  \node[encja] (re) at (0,-3.3) {\naglowek{rezerwacje}
    \nodepart{two} id\_rezerwacji \pk\\ sklep\\ sklep\_zlecajacy\\ klient\\ id\_produktu \fk\\ ilosc\\ status\\ data\_utworzenia};
  \node[encja] (po) at (4.2,-3.3) {\naglowek{pozycje\_zamowien}
    \nodepart{two} id\_zamowienia \pkfk\\ id\_produktu \pkfk\\ ilosc\\ cena\_jednostkowa};
  \node[encja] (za) at (8.4,-3.3) {\naglowek{zamowienia}
    \nodepart{two} id\_zamowienia \pk\\ sklep\\ data\_zamowienia\\ status\\ wartosc};
  \node[encja, dashed] (ra) at (12.6,0) {\naglowek{raporty\_dzienne}
    \nodepart{two} data\_raportu \pk\\ sklep \pk\\ liczba\_zamowien\\ suma\_sprzedazy};

  \draw[rel] (ka.east |- ka.text split) -- (ty.west |- ka.text split)
    node[krt, pos=0.1, above] {1} node[krt, pos=0.85, above] {N};
  \draw[rel] (ty.east |- st.text split) -- (st.west |- st.text split)
    node[krt, pos=0.1, above] {1} node[krt, pos=0.8, above] {0..1};
  \draw[rel] (ty.south) -- (po.north)
    node[krt, pos=0.12, right] {1} node[krt, pos=0.88, right] {N};
  \draw[rel] (ty.south west) -- (re.north east)
    node[krt, pos=0.08, left] {1} node[krt, pos=0.92, above right] {N};
  \draw[rel] (za.west |- po.text split) -- (po.east |- po.text split)
    node[krt, pos=0.1, above] {1} node[krt, pos=0.85, above] {N};
\end{tikzpicture}
\end{center}
{\small Ciemny nagłówek -- tabela fragmentowana (każdy sklep ma własne dane);
jasny -- tabela replikowana z Oracle; linia przerywana -- tabela tylko w SQL Serverze.
Oba sklepy mają ten sam układ tabel.\par}

\subsection{Połączenie Oracle i SQL Servera}
Oba serwery pracują na osobnych laptopach. Połączenie przygotowano w kolejnych krokach,
a każdy z nich sprawdzano przed przejściem do następnego.

\begin{enumerate}
  \item \textbf{Sieć Tailscale.} Oba laptopy zalogowano do tej samej prywatnej sieci
        Tailscale. Każdy dostał stały adres z puli \texttt{100.x.y.z}, a wygasanie
        kluczy wyłączono, żeby połączenie nie przerwało się w trakcie prezentacji.
        Komputery widzą się tak, jakby były w jednej sieci lokalnej, niezależnie
        od tego, gdzie się znajdują.
  \item \textbf{Oracle dostępny z sieci.} Domyślnie listener Oracle nasłuchiwał tylko
        na \texttt{127.0.0.1}. W pliku \texttt{listener.ora} dopisano adres Tailscale,
        dodano regułę zapory dla portu 1521 i zrestartowano listener. Polecenie
        \texttt{lsnrctl status} potwierdziło oba adresy i usługę \texttt{freepdb1}.
  \item \textbf{SQL Server dostępny z sieci.} W SQL Server Configuration Manager
        włączono protokół TCP/IP na porcie 1433, dodano regułę zapory i włączono
        logowanie mieszane (Windows i SQL Server), potrzebne dla loginu aplikacji.
  \item \textbf{Test portów.} Przed uruchomieniem baz sprawdzono samą sieć poleceniem
        \texttt{Test-NetConnection} w obie strony (porty 1521 i 1433). Zwykły
        \texttt{ping} nie jest tu miarodajny, bo zapora często go blokuje.
  \item \textbf{Linked Server.} Na laptopie z SQL Serverem zainstalowano 64-bitowy
        sterownik Oracle Provider for OLE DB (pakiet ODAC) i utworzono Linked Server
        \texttt{ORACLE\_WRO}, wskazujący adres Tailscale Wrocławia i bazę
        \texttt{FREEPDB1}. Wyłączono w nim przenoszenie transakcji na serwer zdalny,
        żeby SQL Server nie próbował tworzyć transakcji rozproszonej (MSDTC).
  \item \textbf{Pierwsza replikacja.} Procedura \texttt{sp\_synchronizuj\_katalog}
        skopiowała kategorie i typy produktów z tymi samymi identyfikatorami co w Oracle
        (w kopii klucze nie są generowane automatycznie). Dopiero potem wstawiono stany
        magazynowe Warszawy, celowo z innymi cenami niż we Wrocławiu.
  \item \textbf{Aplikacja na obu stanowiskach.} Oba laptopy uruchamiają ten sam kod.
        Różnią się tylko plikiem \texttt{.env} (sklep stanowiska, adresy i hasła), który
        nie trafia do repozytorium. Do połączenia z SQL Serverem potrzebny jest sterownik
        ODBC Driver 18.
\end{enumerate}

\subsection{Replikacja katalogu}
Procedura replikacji działa w dwóch etapach. Najpierw pobiera aktualny katalog z Oracle
do tabel tymczasowych, poza transakcją. Następnie w jednej lokalnej transakcji nanosi
zmiany na kopię instrukcją \texttt{MERGE}: dodaje nowe produkty, aktualizuje zmienione,
a produkty usunięte w Oracle oznacza jako nieaktywne zamiast je kasować (mogą być do nich
zamówienia). Jeśli Oracle nie odpowiada, błąd pojawia się już w pierwszym etapie i kopia
zostaje nienaruszona, a Warszawa pracuje na ostatniej pobranej wersji. Każde uruchomienie
zapisuje się w tabeli \texttt{log\_synchronizacji}.

Login, którym aplikacja łączy się z SQL Serverem, może czytać i zapisywać dane sklepu
oraz uruchamiać procedury, ale ma zabronione (\texttt{DENY}) zmiany w kopii katalogu.

\subsection{Sposób uruchamiania i testowania}
Tabele tworzą skrypty z katalogu \texttt{sql}. Aplikację uruchamia się poleceniem
\texttt{streamlit run app/main.py}, a sklep stanowiska ustawia się w pliku \texttt{.env}.
Skrypt \texttt{app/test\_systemu.py} automatycznie sprawdza połączenia, tabele,
fragmentację, replikację, widoki globalne, operacje aplikacji i zachowanie przy
awarii serwera. Skrypt \texttt{app/test\_wspolbieznosc.py} sprawdza, że przy dwóch
jednoczesnych rezerwacjach ostatniej sztuki udaje się tylko jedna.

Scenariusze sprawdzane ręcznie w aplikacji:

{\small
\begin{tabularx}{\textwidth}{@{}>{\raggedright\arraybackslash}p{0.33\textwidth}X@{}}
\toprule
\textbf{Działanie} & \textbf{Oczekiwany wynik} \\
\midrule
Zmiana nazwy produktu w Oracle i synchronizacja & nowa nazwa w kopii w SQL Serverze, polskie znaki poprawne \\
Dodanie produktu w katalogu we Wrocławiu & produkt w Oracle; w Warszawie po pobraniu katalogu \\
Przyjęcie nowego produktu z ceną & produkt pojawia się w stanach sklepu \\
Podgląd produktów w sieci & ceny i ilości z obu sklepów w jednej tabeli \\
Rezerwacja towaru z drugiego sklepu & rezerwacja zapisana w bazie sklepu z towarem i widoczna w obu stanowiskach \\
Sprzedaż & nowe zamówienie we własnej bazie, stan zmniejszony \\
Zatrzymanie jednego serwera lub Tailscale & węzeł oznaczony jako niedostępny, drugi sklep działa normalnie \\
\bottomrule
\end{tabularx}\par}

\douzup{wyniki testów i zrzuty ekranu}

% ============================================
\section{Opis implementacji aplikacji dostępowej}
Aplikacja składa się z czterech plików: \texttt{config.py} (ustawienia),
\texttt{bazy.py} (połączenia z bazami), \texttt{operacje.py} (sprzedaż, rezerwacje,
magazyn, katalog) i \texttt{main.py} (interfejs). Zapytania SQL są pisane raz
i działają w obu bazach. Gdy jedna baza nie odpowiada, aplikacja pokazuje
ostrzeżenie i nadal wyświetla dane z drugiej.

% ============================================
\section{Podsumowanie}
Zrealizowano rozproszoną bazę danych złożoną z Oracle i SQL Servera na dwóch
komputerach, wraz z aplikacją dostępową. Wykorzystano fragmentację poziomą,
replikację, Linked Server, widoki globalne i obsługę awarii. Najważniejszy wniosek:
dane rozmieszczono tak, by prawie każda operacja dotyczyła jednej bazy, bo Oracle
i SQL Server nie mają wspólnych transakcji. Najwięcej pracy wymagało samo połączenie
heterogenicznych serwerów: udostępnienie Oracle w sieci, sterownik OLE DB dla Linked
Servera i zgodność typów danych między Oracle a SQL Serverem. Pomogło sprawdzanie
każdego etapu osobno (najpierw sieć, potem bazy, na końcu replikacja).

\douzup{wnioski z pomiaru czasu zapytań lokalnych i przez Linked Server}

% ============================================
\begin{thebibliography}{9}\small
\bibitem{ozsu} M. T. Özsu, P. Valduriez, \emph{Principles of Distributed Database Systems}, wyd. 4, Springer, 2020.
\bibitem{wyklad} \douzup{materiały do wykładu}
\bibitem{mslinked} Microsoft, \emph{Linked Servers (Database Engine)},
  \url{https://learn.microsoft.com/sql/relational-databases/linked-servers/linked-servers-database-engine}
\bibitem{oradoc} Oracle, \emph{Oracle Database Documentation}, \url{https://docs.oracle.com/en/database/}
\end{thebibliography}

\end{document}
